

\section{Introduction}

Automatic evaluation metrics are a cornerstone of summarization research,
enabling rapid and inexpensive iteration during model development. Yet the two
dominant paradigms, reference-based and reference-free evaluation,
each carry a fundamental limitation that constrains their practical utility.

Given recent success building high-performance Large Language Models (LLMs), researchers have also explored using LLMs to evaluate the generated output of other LLMs, a setting known as LLM-as-Judge. LLM-as-judge evaluation~\cite{zheng2023judging,liu2023geval}
leverages the reasoning capabilities of LLMs to score generated summaries,
achieving high correlation with human judgment across a range of text
generation tasks. Crucially, LLM judges can operate in a
\textit{reference-free} setting~\cite{Deutsch2022LimitationsReferenceFree},
assessing summary quality without any gold reference. This completely removes the
annotation bottleneck. However, reference-free LLM judges
lack an external anchor for their assessments, and prior work shows that they
perform substantially better when a reference is
provided~\cite{zheng2023judging,kim2023prometheus}.

This creates a practical dilemma for researchers and developers. Gold
references yield better evaluations, but collecting them is expensive.
Reference-free evaluation is cheap but weaker. In many real-world settings,
this tradeoff is particularly acute: a researcher evaluating a new
domain-adapted summarization model may have access to \textit{related}
summaries, such as abstracts of scientific articles, or a handful of
human-written examples from a pilot annotation, but not a gold reference
for every test document. Under current paradigms, this approximate information is not leveraged. 




\textbf{We propose reference-adjusted evaluation}, a new paradigm that
occupies the practical middle ground between the two current standards. Rather than requiring exact gold references or discarding all reference information
entirely, reference-adjusted evaluation leverages \textit{related but
non-exact} references: summaries that share either the desired content or the
target style of the ideal output, but not necessarily both. For example, when
evaluating lay summaries of scientific articles, document abstracts are
naturally available. They share the desired informational content but are
written for an expert audience rather than a lay one. Similarly, a small set
of gold lay summaries from a pilot study shares the target style but covers
different documents than the test set. In both cases, the related summary
provides a useful reference point for the LLM judge without requiring a
full annotation effort. 

We introduce the reference-adjusted evaluation paradigm and demonstrate that
it improves LLM-as-judge performance over the reference-free setting across
three open-source models and four evaluation axes, using cheaply acquired
machine-generated summaries as approximate references. We highlight a practical path forward for researchers who need robust evaluation without the cost of full gold annotation. We visualize the trade-offs of the 3 evaluation paradigms in~\Cref{fig:ref-adj-vis}. Reference-based evaluation has high correlation with human judgment, but lacks scalability and ease of deployment, while requiring high annotation cost. Reference-free evaluation is scalable, easy to deploy, and requires no annotation, but suffers from lower correlation with human judgment. Reference-adjusted evaluation bridges this gap by improving ease of deployment, scalability, low to no annotation cost, and higher correlation with human judgment. 